\documentclass[journal]{IEEEtran}

%% Use Latin Modern for the entire family (text + tt) so the local TeX
%% install does not try to call Metafont for pcr Courier metrics that
%% live in the (separately distributed) psnfss bundle.
\usepackage[T1]{fontenc}
\usepackage{lmodern}
\renewcommand*{\ttdefault}{lmtt}

\usepackage{cite}
\usepackage{amsmath,amssymb,amsfonts}
\usepackage{graphicx}
\usepackage{textcomp}
\usepackage{xcolor}
\usepackage{booktabs}
\usepackage{array}
\usepackage[colorlinks=true,
            linkcolor=blue,
            citecolor=blue,
            urlcolor=blue]{hyperref}

%% Fix bad-break tolerance for IEEE journal two-column.
\setlength{\emergencystretch}{2em}
\widowpenalty=10000
\clubpenalty=10000

%% Allow ragged-right + a much larger emergency stretch inside the
%% bibliography so long urls / paths in the citations wrap rather than
%% overrun column boundaries.
\let\oldbibliography\bibliography
\renewcommand{\bibliography}[1]{%
  \begingroup
  \sloppy
  \raggedright
  \setlength{\emergencystretch}{6em}%
  \oldbibliography{#1}%
  \endgroup
}

\title{RayaChain: An Omnichain Canary Layer for
       Cross-Chain Risk and Credit Signal Propagation in
       Decentralized Agentic Finance}

%% Author block --- IEEEtran journal style.
%% Author 2 details (surname / email / affiliation / ORCID / IEEE membership)
%% are AWAITING CEO confirmation. Placeholders in the \thanks below must be
%% filled before submission. Do not submit with placeholders.
\author{Basel~Magableh,~\IEEEmembership{Member,~IEEE,}
        and~Ahod~[surname~TBD]%
\thanks{B. Magableh is with the School of Computer Science, Technological
University Dublin, Dublin, Ireland (e-mail: basel.magableh@tudublin.ie).
ORCID: 0000-0003-2337-637X.}%
\thanks{A. [surname TBD] is with [affiliation TBD] (e-mail: [email TBD]).}%
\thanks{Manuscript received May~7,~2026.}}

\markboth{IEEE Transactions on [venue~TBD],~Vol.~XX, No.~X, May~2026}%
{Magableh \MakeLowercase{\textit{et al.}}: RayaChain Omnichain Canary Layer}

\IEEEpubid{0000--0000/00\$00.00~\copyright~2026 IEEE}

\begin{document}

\maketitle

\begin{abstract}
Cross-chain interoperability has enabled the seamless movement of value
across blockchain ecosystems, yet risk and credit intelligence remain
fragmented, chain-local, and reactive. Existing interoperability
protocols transport value and arbitrary messages but do not propagate
\emph{evaluated risk context} or \emph{portable credit intelligence}
across heterogeneous networks. This paper introduces \textbf{RayaChain},
an omnichain \emph{canary layer} for cross-chain risk and credit signal
propagation in decentralized agentic finance, in which autonomous
agents---rather than human traders alone---consume the propagated
context to drive routing, lending, and liquidity decisions. RayaChain is
built on LayerZero OFT messaging and composes signal generation,
semantic encoding, cross-chain transmission, and signal-consumption
layers into a single primitive. We formalise the canary-layer
abstraction, define wallet-, protocol-, and market-level signals, and
evaluate the system across exploit-propagation, credit-portability, and
liquidity-stress scenarios. We further argue that cross-chain risk
intelligence becomes verifiable when paired with a Bittensor-style
decentralised validation subnet, and report orchestration evidence
($5{,}097$ successful rounds, zero authentication failures) from a
testnet shadow deployment. The contribution is a missing infrastructure
primitive for the decentralized agentic-finance stack: an
interoperability layer for \emph{intelligence}, not value.
\end{abstract}

\begin{IEEEkeywords}
Blockchain, decentralized finance, cross-chain interoperability, risk
intelligence, wallet scoring, LayerZero, omnichain systems, Bittensor,
agentic finance, constitutional AI.
\end{IEEEkeywords}

%% =====================================================================
\section{Introduction}
%% =====================================================================

\IEEEPARstart{D}{ecentralized} finance (DeFi) has evolved into a
heterogeneous, multi-chain ecosystem in which value flows freely across
networks~\cite{BelchiorVGC22,00330023}. The next architectural shift,
already visible in production, replaces the human trader at the
edge of that system with \emph{autonomous agents}: LLM- or policy-driven
processes that observe on-chain and off-chain context, select an action
under role-bound constitutional
constraints~\cite{bai2022constitutional}, and execute against
cross-chain venues without a human in the loop for every step. We refer
to this regime as \emph{decentralized agentic finance}.

Despite the maturity of cross-chain \emph{value} transfer, risk
intelligence remains siloed within individual chains. Wallet
reputation, behavioural risk, protocol exposure, and route health do
not travel with assets~\cite{HasanBB23}. The blind spots this creates
are not merely operational: an autonomous agent acting on incomplete
cross-chain context will, by construction, route capital into
exploit-active venues, refuse credit to wallets whose entire repayment
history lives on a different chain, and miss liquidity stress that is
visible one bridge away.

This paper introduces the concept of a \textit{canary layer}: a
distributed early-warning system for risk propagation across chains.
We present \textbf{RayaChain} as an implementation of the concept,
positioned within the emerging stack of decentralized agentic finance
where autonomous agents take liquidity-, routing-, and credit-allocation
decisions on behalf of human and protocol principals. RayaChain is
built on LayerZero's omnichain messaging
substrate~\cite{abs-2110-13871} and is designed to interoperate with a
decentralised validation subnet for verifiable risk
intelligence~\cite{rao2021bittensor,steeves2022incentivizing,opentensor2024yuma}.

\paragraph*{Contributions.}
This paper contributes:
\begin{enumerate}
\item A formal definition of the \emph{canary layer} as a missing
      infrastructure primitive for decentralized agentic finance: an
      interoperability layer for evaluated intelligence, separate from
      and complementary to value- and message-transport layers.
\item A cross-chain risk- and credit-signal propagation model with a
      three-tier signal taxonomy (wallet, protocol, market) and a
      semantic encoding distinct from raw bridge messaging.
\item A \emph{verifiable risk intelligence} layer that pairs the canary
      layer with a Bittensor-style validation subnet, with a formal
      validator-loss decomposition (Brier accuracy + paired-query
      inconsistency + binned miscalibration) and an honest split between
      orchestration evidence (already in hand: $5{,}097$ successful
      rounds, zero authentication failures over a 57-hour testnet
      shadow soak~\cite{omnirisk2026soak}) and economic verification at
      scale (deferred).
\item Three use-case scenarios (exploit propagation, credit
      portability, liquidity stress) and a metric set (detection
      latency, signal accuracy, cross-chain coverage) that ground the
      empirical agenda.
\end{enumerate}

\paragraph*{Roadmap.}
Section~\ref{sec:bg} states the background and the problem.
Section~\ref{sec:rw} reviews interoperability, analytics, wallet
scoring, and decentralised validation networks.
Section~\ref{sec:arch} specifies the RayaChain architecture and signal
taxonomy. Section~\ref{sec:vri} states the verifiable
risk-intelligence layer formally. Sections~\ref{sec:uc}--\ref{sec:eval}
ground the architecture in three use cases and the experimental design
that operationalises them. Section~\ref{sec:disc} states benefits,
limitations, and future work, and Section~\ref{sec:concl} closes.

%% =====================================================================
\section{Background and Motivation}\label{sec:bg}
%% =====================================================================

\subsection{Cross-Chain Value vs.\ Cross-Chain Intelligence}

Blockchain technology has emerged as a decentralised paradigm for
secure, peer-to-peer value exchange without reliance on centralised
authorities. A blockchain operates as a distributed ledger that records
and verifies transactions through consensus, ensuring immutability,
transparency, and trust among untrusted participants~\cite{MukhopadhyaySHO16,
00330023}. Cryptocurrencies and decentralised applications (DApps)
leverage these properties to enable financial interactions, asset
ownership, and programmable logic via smart contracts~\cite{KhanAAA21a}.

The ecosystem has expanded into a fragmented landscape of heterogeneous
networks with distinct protocols, consensus models, and execution
environments. This fragmentation limits interoperability: blockchains
are inherently isolated and cannot natively communicate or exchange
information across network boundaries~\cite{BelchiorVGC22}. Users and
applications are therefore constrained within individual ecosystems,
producing inefficiencies in asset utilisation and system-wide
coordination.

Cross-chain interoperability solutions, particularly cross-chain
bridges, address this gap by transferring assets and data between
disparate networks. Such bridges enable token mobility and unlock
multi-chain DeFi~\cite{ZhangZZL24}. They operate fundamentally at the
asset layer and do not propagate the contextual or behavioural
information attached to those assets. Existing validation mechanisms
verify the correctness of transactions but do not assess the
\emph{intent} or behavioural patterns of participating
entities~\cite{bhowmik2021}. As a result, transactions deemed valid by
consensus may still encode malicious or fraudulent activity, exposing
a structural absence of intelligence-driven risk assessment.

Cross-chain bridges introduce significant security and architectural
challenges. Their design---verification mechanisms, relayers, and smart
contracts---creates expanded attack surfaces that have been exploited
repeatedly, with cumulative losses in the multi-billion-dollar
range~\cite{MishraP25,LiQXZZWX25,augusto2024sok}. The interoperability
infrastructure remains immature, with unresolved trust assumptions,
verification models, and systemic
vulnerabilities~\cite{MishraP25,werner2022sok}. Existing cross-chain
systems therefore enable the mobility of \emph{value} but not the
mobility of \emph{intelligence}: risk signals, behavioural history, and
trust assessments do not travel with the assets they describe. This
asymmetry is the architectural gap this paper closes.

\subsection{Problem Statement}

Despite the rapid growth of decentralised finance and multi-chain
ecosystems, current blockchain infrastructures lack a unified mechanism
for sharing risk, trust, and behavioural context across chains. The
problem decomposes into five dimensions.

\subsubsection*{Signal fragmentation}
Wallet reputation, protocol exposure, and anomalous behaviour are
visible only on the chain where they originate. Reputation systems
that are essential for evaluating trustworthiness in decentralised
environments aggregate user interactions and feedback to derive trust
scores~\cite{HasanBB23}, but are platform-specific: a user's reputation
on one chain does not transfer to another.

\subsubsection*{Reliability and manipulation}
Existing reputation frameworks suffer from limitations in reliability,
portability, and resistance to manipulation, including false feedback,
reputation fraud, Sybil attacks, and the absence of standardised
aggregation mechanisms~\cite{AlmasoudHH20}. Most systems do not treat
reputation as a transferable, composable digital asset, preventing
integration into broader decentralised ecosystems. Traditional
cryptography and access control address external threats but do not
capture the dynamic behaviour of internal entities, motivating trust
management systems that evaluate behavioural patterns and interaction
histories~\cite{ArshadKAKYZ23}.

\subsubsection*{Latency}
The absence of real-time, cross-chain risk intelligence delays
detection of malicious behaviour~\cite{MishraP25,NahaZ24}. Blockchain
transactions are transparent, yet the lack of a shared intelligence
layer means risk signals---suspicious transaction patterns, wallet
behaviours, or liquidity anomalies---are confined within individual
chains. Malicious actors exploit this fragmentation by moving assets
rapidly across chains before chain-local detection can
respond~\cite{DBLP:journals/spe/DelgadovonEitzenRMF25,DudaniBRM23}.

\subsubsection*{Systemic risk in agentic regimes}
Reliance on cross-chain infrastructure amplifies systemic
risk~\cite{LiQXZZWX25,ZhangZZL24}. As assets move seamlessly across
chains, vulnerabilities in one ecosystem propagate indirectly to
others. The agentic regime sharpens the consequence: an autonomous
agent that allocates capital under a constitutional
contract~\cite{bai2022constitutional} can only be safe if its action
surface is constrained by the same risk context human operators would
demand. Without a cross-chain intelligence layer, the agent acts on
chain-local information and cascades the exposure.

\subsubsection*{Credit portability}
A wallet may exhibit a strong repayment or behavioural history on one
chain and appear as a blank identity on another, preventing the
emergence of portable on-chain credit and risk-aware
under-collateralised finance.

Collectively, these dimensions reveal that while value moves freely
across chains, intelligence does not. Closing the gap requires an
interoperable intelligence layer capable of aggregating, analysing, and
propagating risk and trust signals across heterogeneous networks in a
form that autonomous agents can consume.

\subsection{Research Questions and Hypotheses}\label{sec:rq}

We re-cast the problem as five research questions and four
hypotheses, framed for the agentic regime.

\paragraph*{Main research question}
\textbf{RQ1.} How can cross-chain systems propagate evaluated wallet-,
protocol-, and market-level risk signals in a timely and trustworthy
way to support proactive decision-making by autonomous agents in
decentralized finance?

\paragraph*{Supporting questions}
\textbf{RQ2.} What classes of on-chain behaviour are most predictive as
early-warning signals when aggregated across multiple chains?

\textbf{RQ3.} How should cross-chain risk signals be encoded,
transported, and consumed so that they remain interpretable and useful
to downstream protocols and agentic policies?

\textbf{RQ4.} What trade-offs emerge between signal speed, reliability,
granularity, and false-positive risk in an omnichain risk-distribution
architecture?

\textbf{RQ5.} To what extent can portable cross-chain behavioural
history support a practical notion of decentralised creditworthiness
under constitutional agentic constraints?

\paragraph*{Hypotheses}
\textbf{H1.} Cross-chain propagation of evaluated risk signals reduces
detection latency relative to chain-local monitoring alone.

\textbf{H2.} Multi-chain behavioural aggregation produces more
predictive early-warning indicators than single-chain observations.

\textbf{H3.} Portable risk and credit context improves downstream
decision quality for lending, exchange screening, and liquidity
management---and, in the agentic regime, narrows the action surface of
constitutionally-constrained agents to safer policies.

\textbf{H4.} A semantically-aware signal layer provides better
operational utility than generic cross-chain messaging alone, because
agentic consumers need decision-useful surfaces rather than raw bytes.

%% =====================================================================
\section{Related Work}\label{sec:rw}
%% =====================================================================

\subsection{Cross-Chain Interoperability Protocols}

The expansion of blockchain ecosystems has produced multiple
interoperability solutions for communication across heterogeneous
networks~\cite{BelchiorVGC22}. LayerZero identifies fragmentation as a
direct consequence of the proliferation of independent blockchains, in
which each chain operates in isolation and forces users to split
liquidity, applications, and state across disconnected
ecosystems~\cite{abs-2110-13871}. Smart-contract-based approaches
enable programmable coordination between chains, allowing decentralised
applications to interact beyond a single network~\cite{KhanAAA21a}.
LayerZero extends this paradigm with a trustless omnichain
interoperability protocol that functions as a low-level messaging
layer, enabling direct cross-chain communication without intermediary
chains or custodial entities~\cite{abs-2110-13871}. Practical
instantiations include LayerZero's Omnichain Fungible Token (OFT) and
Chainlink's Cross-Chain Interoperability Protocol (CCIP).

These protocols remain fundamentally focused on \textit{transport}.
They move assets and data but do not account for the semantic meaning
or contextual interpretation of that data. They lack mechanisms for
propagating risk signals, behavioural insights, or trust-related
information. Value moves seamlessly between ecosystems; the
intelligence required to interpret that movement remains confined to
its origin chain. Security challenges compound the
limitation: cross-chain bridges introduce significant attack
surfaces~\cite{ZhangZZL24,LiQXZZWX25,augusto2024sok}, and the
ecosystem still lacks mature public mechanisms for collecting,
organising, and visualising cross-chain operations at
scale~\cite{SantosAVC25}, constraining security analysis and timely
understanding of bridge activity.

\subsection{Blockchain Analytics Platforms}

Blockchain analytics platforms have emerged as critical tools for
understanding on-chain behaviour and detecting anomalies. Platforms
such as Nansen and Arkham Intelligence apply transaction-graph
analysis, clustering, and labelled datasets to identify entities, trace
fund flows, and detect suspicious patterns. Their foundations lie in
graph-based representations of blockchain
systems~\cite{MukhopadhyaySHO16,SerenaFD22}, extended to multi-layered
graph systems that incorporate distinct interaction types---money
flows, contract creation, contract invocation---to enable fine-grained
behavioural and security analysis~\cite{ChenLZCLLLZ20}. Empirical
studies confirm that cryptocurrency transaction networks exhibit
power-law degree distributions and non-assortative behaviour, making
them suitable for large-scale behavioural and network-level
analysis~\cite{MotamedB19}. Recent work integrates security and
forensic perspectives with machine learning and behavioural clustering
for fraud and illicit-activity
detection~\cite{MishraP25,FanLZJB20,electronics13173568}; existing
forensic approaches still lack integrated support for real-time
analytics and graph-based storage critical to modern
investigations~\cite{NahaZ24,MasudHSAAAY21}.

These platforms exhibit three structural limitations. First, they
operate primarily as \textbf{centralised intelligence layers}, relying
on proprietary data pipelines and off-chain processing. Second, their
analysis is largely \textbf{chain-specific}, interpreting behavioural
patterns within a single blockchain. Third, and most critically, they
lack mechanisms for \textbf{real-time intelligence propagation}: risk
signals identified on one chain are not communicated to other
ecosystems, limiting effectiveness in multi-chain environments and
making the platforms unsuitable as direct inputs to autonomous
cross-chain agents.

\subsection{Wallet Scoring and Trust Systems}

Reputation and trust modelling quantify the reliability of blockchain
participants from historical interactions, enabling more informed
decision-making in trustless environments. Traditional reputation
systems aggregate user feedback and behavioural data~\cite{HasanBB23};
blockchain-based implementations enhance transparency and
tamper-resistance, making them suitable for decentralised
ecosystems. Smart-contract-based reputation systems automate trust
management and policy enforcement~\cite{AlmasoudHH20}; a key residual
limitation is the lack of interoperability across reputation
systems---reputation values remain bound to a specific platform.

Machine-learning-based wallet scoring detects fraudulent or high-risk
behaviour by analysing transaction patterns, temporal features, and
network structures~\cite{MishraP25,palaiokrassas2023multichain}.
Multichain DeFi fraud detection across 12 chains and 23 protocols has
been studied empirically~\cite{palaiokrassas2023multichain}, but, as
with analytics platforms, models are typically trained centrally on
isolated datasets and do not propagate the resulting scores back to
the chains where the wallet acts.

\subsection{Decentralised Validation Networks}\label{sec:rw:dvn}

A separate strand of recent work asks how risk-relevant outputs
themselves can be validated in a decentralised, economically incentivised
manner. The canonical instantiation is Bittensor: a peer-to-peer
intelligence market in which a heterogeneous population of \emph{miners}
submits outputs for a task, and \emph{validators} score the miners
against realised
outcomes~\cite{rao2021bittensor,steeves2022incentivizing}. The Yuma
Consensus rule clips outlier validator weight vectors against a
stake-weighted threshold and emits stake-aligned
rewards~\cite{opentensor2024yuma}, producing a coalition-resistant
ranking under stake fractions below a known
bound~\cite{steeves2022incentivizing}. The mechanism-design lineage is
the prediction-market literature: Hanson's logarithmic market scoring
rule is the canonical primitive that aggregates self-interested
predictors into a calibrated forecast~\cite{hanson2003combinatorial},
and Taoshi's Subnet 8 (Proprietary Trading Network) is the closest
live-2026 precedent for a financial-prediction
subnet~\cite{taoshi2024ptn}. A companion focused
manuscript~\cite{magableh2026agenticiov} composes a five-engine
architecture in which a Bittensor verification subnet decentralises a
prediction engine and economically scores its outputs against
realised cross-chain events; we adopt the same framing here, but apply
it specifically to the canary-layer signal substrate.

The relevance of decentralised validation networks to the canary layer
is twofold. First, an autonomous agent that consumes propagated risk
signals must trust the signal source; a Bittensor-scored signal is
\emph{economically} backed rather than reputationally backed. Second,
the validator-loss decomposition formalised in
Section~\ref{sec:vri:formal} is falsifiable, so the canary layer
inherits a measurable verification surface rather than a black-box
trust claim.

\subsection{Research Gap}

Existing research provides strong foundations: interoperability
protocols enable the movement of value and data; analytics platforms
provide deep insights into on-chain behaviour; reputation systems offer
mechanisms for evaluating trust; and decentralised validation networks
make output quality economically scoreable. Forensic and analytical
frameworks for blockchain remain incomplete, with several investigation
phases either underexplored or
unaddressed~\cite{MasudHSAAAY21,DudaniBRM23}.

The gap is specific. Current systems move \textit{value} and, in some
cases, arbitrary \textit{data}, but they do not move
\textit{intelligence}. Risk signals, behavioural patterns, and credit
information remain fragmented across chains, even as assets and users
operate in increasingly interconnected environments. Worse, the next
generation of consumers of these signals is no longer human: it is
autonomous, constitutionally-constrained agents whose action surface is
defined by the context they can actually
see~\cite{bai2022constitutional}. The result is systemic blind spots in
cross-chain agentic decision-making, not merely in cross-chain trader
decision-making.

This paper addresses the gap by introducing an interoperable
intelligence layer that observes and analyses on-chain behaviour and
propagates structured, actionable, agent-consumable insights across
blockchain ecosystems in real time. Pairing the layer with a
decentralised validation subnet makes the propagated signals
economically verifiable, not merely operationally useful.

%% =====================================================================
\section{RayaChain Architecture}\label{sec:arch}
%% =====================================================================

\subsection{System Overview}

RayaChain is an omnichain signal layer built on LayerZero OFT
infrastructure~\cite{abs-2110-13871}. Where LayerZero transports value
and arbitrary messages between chains, RayaChain transports
\emph{evaluated context}: structured, semantically-typed risk and
credit signals carrying a confidence and a window of validity. The
layer is positioned beneath agentic policy engines and above existing
interoperability primitives. Figure-conceptually, an originating chain
emits a signal; a relayer carries the signal and its provenance to a
target chain; and a consuming protocol or agent reads the signal under
a typed schema before deciding to act.

\subsection{Core Components}

RayaChain decomposes into four components.

\subsubsection*{Signal generation (OmniRisk engine)}
A signal-generation engine ingests on-chain transactions, wallet
behaviour logs, liquidity-pool state, bridge flow, and off-chain
context (news, social streams, grey-literature feeds for emerging
markets) and emits a typed signal of the form
\[
s = (\text{chain}_o, \text{subject}, \text{type}, v, c, w, \text{prov}),
\]
where $\text{chain}_o$ is the origin chain, $\text{subject}$ is the
wallet, protocol, or market identifier, $\text{type}$ is one of the
classes of Section~\ref{sec:arch:taxonomy}, $v$ is the signal value,
$c \in [0,1]$ is a confidence, $w$ is a window of validity, and
$\text{prov}$ is provenance metadata.

\subsubsection*{Signal encoding}
Signals are serialised against a versioned schema with a fixed byte
budget and a deterministic canonical form so that two relayers agreeing
on the same signal produce byte-identical payloads. Encoding decouples
the semantics of a signal from the wire format used by the underlying
transport, permitting the canary layer to ride on any messaging
substrate---LayerZero OFT in this paper---without re-deriving the
type system.

\subsubsection*{Cross-chain transmission}
Encoded signals are transmitted across chains using LayerZero's
ultra-light-node messaging protocol~\cite{abs-2110-13871}. The
canary-layer signal type is a strict subset of LayerZero's payload
contract; we treat the underlying transport as a substitutable bearer.

\subsubsection*{Signal-consumption layer}
On the destination chain, a signal-consumption library exposes a
typed read API to consuming protocols and agents. The consumer can
filter by signal type, freshness window, and minimum confidence, and
receive a typed object rather than raw bytes. This is the surface that
distinguishes the canary layer from LayerZero alone: the consumer never
parses the wire format directly, and a downstream agent can be
constitutionally constrained to only act on signals of a declared
type and confidence band.

\subsection{Schema Excerpt}

The wire-level schema is intentionally narrow. A canonical
specification of the wallet-level signal type, in pseudo-IDL form,
is:
\begin{quote}\small
\texttt{struct WalletSignal \{ chainId origin; bytes20 wallet;
WalletSignalType type; int16 value\_bps; uint8 confidence\_q7; uint64
window\_start; uint64 window\_end; bytes32 prov\_root; \}}
\end{quote}
\noindent
where \texttt{value\_bps} is a fixed-point value in basis points,
\texttt{confidence\_q7} is a Q0.7 fixed-point confidence in $[0,1]$,
$[\texttt{window\_start}, \texttt{window\_end}]$ defines the validity
window, and \texttt{prov\_root} is a Merkle root over the underlying
evidence set. Equivalent structures are defined for the protocol- and
market-level signal classes; the \texttt{prov\_root} field is the hook
that lets a verifying consumer chase signal provenance back to an
attestation set without re-deriving the signal value itself.

\subsection{Canary Layer Definition}

\textbf{Definition (Canary Layer).} A canary layer is an
\emph{early-warning signal distribution system} for risk and credit
context across blockchain ecosystems, with three hard properties:
(i)~\emph{semantic typing}---signals are typed against a shared schema
rather than carried as opaque bytes;
(ii)~\emph{evaluated context}---signals carry a value, a confidence, and
a window of validity, not merely a transaction record; and
(iii)~\emph{cross-chain propagation}---signals are transported across
heterogeneous networks under a single addressable namespace. A canary
layer is therefore distinct from a value-transfer bridge (which moves
assets), an arbitrary-message bridge (which moves bytes), and a
chain-local analytics platform (which observes but does not propagate).

\subsection{Signal Taxonomy}\label{sec:arch:taxonomy}

We define three signal classes, each with a distinct generator, decay
window, and downstream consumer.

\subsubsection{Wallet-level signals}
Behavioural anomalies attached to a wallet identifier: sudden
diversification, interaction with known exploit contracts, liquidity
withdrawal cascades, or repayment-failure history. Generators include
graph-analytic
features~\cite{ChenLZCLLLZ20,SerenaFD22,MotamedB19}, behavioural
clustering~\cite{FanLZJB20}, and cross-chain reputation
aggregation~\cite{HasanBB23,AlmasoudHH20}. Consumers include lending
protocols (collateral haircuts), exchange-screening agents, and
agent-task-management layers that gate which wallets an autonomous
trader role may interact with.

\subsubsection{Protocol-level signals}
Liquidity-stress and bridge-anomaly indicators attached to a protocol
or pool identifier: utilisation spikes, oracle divergence, bridge-flow
asymmetries, and unusual TVL contraction. Generators draw on bridge
forensics and analytics literature~\cite{MishraP25,LiQXZZWX25,
ZhangZZL24,augusto2024sok} and on AMM-level liquidity studies that
characterise how the same logical asset prices differently across
venues~\cite{gogol2024liquidity}. Consumers include router agents,
liquidity-management agents, and protocol-level circuit breakers.

\subsubsection{Market-level signals}
Systemic risk patterns attached to a market or correlated cluster of
assets: liquidity-flight indicators, narrative shocks
detectable as fusion-divergence between text-stream and on-chain-flow
sentiment~\cite{cong2022crypto}, and sandwiching or MEV-active
windows~\cite{daian2020flash}. Consumers are higher-order agents that
must rebalance or pause when a systemic regime change is detected
within the agent's constitutional action surface.

%% =====================================================================
\section{Verifiable Risk Intelligence}\label{sec:vri}
%% =====================================================================

A propagated signal is only as good as the trust that downstream
agents can place in it. The canary layer of Section~\ref{sec:arch}
provides typed propagation; it does not, by itself, provide a
verifiable backing for the values it carries. This section adds that
backing by composing the layer with a decentralised validation
subnet, and states the specification formally.

\subsection{Verification by Decentralised Subnet}

We exposed the signal-generation engine of Section~\ref{sec:arch} as
a Bittensor subnet, in which a heterogeneous population of miners
emits forecasts of subject-level risk events and validators score the
forecasts against realised outcomes under the Yuma Consensus
rule~\cite{rao2021bittensor,steeves2022incentivizing,opentensor2024yuma}.
The mechanism-design lineage is Hanson's market scoring
rule~\cite{hanson2003combinatorial}, and the closest live-2026
precedent for a Bittensor financial-prediction subnet is Taoshi's
Subnet 8~\cite{taoshi2024ptn}.

\subsection{Formal Validator-Loss Specification}\label{sec:vri:formal}

We formalise the validator-side loss; the specification is reproduced
from the companion focused manuscript~\cite{magableh2026agenticiov}
and is a design proposal that we explicitly do not yet measure
against a live multi-miner metagraph.

\paragraph*{Miner output.}
A miner $i$ emits, for each query, a pair
\begin{equation}
m_i = (p_i, c_i),
\end{equation}
where $p_i \in [0,1]$ is the predicted probability of the realised
event over a scoring window $\Delta$, and $c_i \in [0,1]$ is the miner's
self-stated calibration confidence in $p_i$.

\paragraph*{Material chain-state event.}
The realised-event indicator $y \in \{0,1\}$ is set to $1$ if any of
four event classes occurs within $\Delta$ of the prediction emission:
$E_1$, a route-level liquidity contraction beyond a configured
threshold; $E_2$, a price drop beyond a configured threshold over a
configured sub-window; $E_3$, a bridge or oracle anomaly registered by
a cross-source redundancy check; $E_4$, a governance state change
recorded on-chain. The thresholds and the union over $E_1$--$E_4$
together define the resolver.

\paragraph*{Validator loss.}
The validator-side loss for miner $i$ is the convex combination
\begin{align}
L_i ={}& \alpha\,\mathrm{Brier}(p_i,y)
       + \beta\,\mathrm{Inconsistency}(m_i) \notag\\
      &+ \gamma\,\mathrm{Calibration}(c_i),
\end{align}
with $\alpha+\beta+\gamma = 1$.
The Brier component $\mathrm{Brier}(p,y) = (p-y)^2$ penalises both
wrong predictions and overconfidence. The inconsistency component
penalises drift across paired queries that share the same chain-state
context (no material change between query times). The calibration
component is the per-bin gap between stated confidence $c_i$ and
realised accuracy over the past $W$ paired $(c_i,y)$ samples.

\paragraph*{Defaults.}
We propose $\alpha = 0.6$, $\beta = 0.2$, $\gamma = 0.2$ as the design
starting point. The Brier component carries the largest weight because
it is the only component grounded directly in realised outcomes; the
inconsistency and calibration components are intended to be hard to
game without also producing accurate forecasts.

\paragraph*{Killable-uplift property.}
The verification subnet is a \emph{killable uplift} to the
signal-generation engine: if the empirical distribution of $L_i$ across
active miners has too high a variance over a fixed validation window,
the operator can disable the subnet path without losing the production
signal surface. This bounds the worst-case verification regression to
the signal-generation engine's standalone calibration error.

\subsection{Threat Model and Trust Assumptions}\label{sec:vri:threat}

The verifiable canary layer is designed to tolerate a small but
load-bearing set of adversaries. Table~\ref{tab:threat-model} states
each as a row over capability, architectural defence, and the residual
surface the architecture does not cover. The model is adapted from the
companion focused
manuscript~\cite{magableh2026agenticiov} and tightened for the
canary-layer signal substrate.

\begin{table*}[t]
\centering
\footnotesize
\caption{Threat model for the verifiable canary layer. Each row gives
the adversary's assumed capability, the architectural defence, and the
residual surface the architecture does not cover.}
\label{tab:threat-model}
\renewcommand{\arraystretch}{1.15}
\begin{tabular}{@{}p{0.16\linewidth} p{0.27\linewidth} p{0.27\linewidth} p{0.22\linewidth}@{}}
\toprule
\textbf{Adversary} & \textbf{Capability} & \textbf{Defence} &
\textbf{Residual risk} \\
\midrule
Malicious miner & Arbitrary forecast emission, adaptive to validator
queries & Three-component validator loss
(Section~\ref{sec:vri:formal}); killable-uplift property of the subnet
& Miner that learns the query distribution may game the inconsistency
component \\
Colluding validators & Coordinated weight submission across stake
fraction $f$ & Clipped weights and stake distribution under the Yuma
rule~\cite{opentensor2024yuma} & A $\geq 50\%$ stake coalition can
still misrank; mitigated by stake-distribution governance \\
Sentiment poisoning & Coordinated text injection across news, social,
and grey-literature feeds & Provenance-tracked ingestion;
deterministic lexicon; fusion-divergence detector vs.\ on-chain-only
baseline & Attacker that also influences on-chain flow (e.g.\ wash
trades) defeats the divergence detector~\cite{cong2022crypto} \\
Oracle / route manipulation & Control over a single oracle or venue
quote, including sandwiching~\cite{daian2020flash} & Cross-source
redundancy in the signal-generation engine; divergence flags cohort
\textsc{degraded} & Correlated manipulation across all sources is
not detected \\
Replay attacker & Re-emission of past signed signals & Window-aligned
canonical timestamps; out-of-window signals rejected & Same-window
replay between adjacent miners is indistinguishable from genuine
convergence \\
Bridge-substrate compromise & Exploit of the underlying LayerZero
substrate~\cite{ZhangZZL24,augusto2024sok} & Killable-substrate
property: the canary layer can be re-pointed to an alternative
transport without re-deriving signal types & Substrate compromise
during the re-point window is unmitigated \\
\bottomrule
\end{tabular}
\end{table*}

\subsection{Validation Status of the Architecture}

Table~\ref{tab:validation-status} states which components of the
architecture are measured in this paper, which are partially validated,
and which are stated architecturally but not yet measured. The table
preempts the reading of the paper as a uniformly empirical artefact:
only the orchestration path is measured at this revision.

\begin{table*}[t]
\centering
\small
\caption{Validation status of the canary-layer components. \textsc{yes}
denotes a live empirical measurement reported in this paper;
\textsc{shadow only} denotes a testnet run not yet against a real
metagraph; \textsc{architectural} denotes a component stated formally
but not yet measured.}
\label{tab:validation-status}
\renewcommand{\arraystretch}{1.15}
\begin{tabular}{@{}>{\raggedright\arraybackslash}p{0.34\linewidth} >{\raggedright\arraybackslash}p{0.30\linewidth} >{\raggedright\arraybackslash}p{0.30\linewidth}@{}}
\toprule
\textbf{Component} & \textbf{Status} & \textbf{Evidence} \\
\midrule
Signal generation engine & \textsc{yes} (in production on the
underlying eight-chain deployment) & Companion
manuscript~\cite{magableh2026agenticiov} \\
Signal encoding (typed schema) & \textsc{architectural} & This paper,
Sec.~\ref{sec:arch} \\
Cross-chain transmission (LayerZero) & \textsc{architectural} & This
paper, Sec.~\ref{sec:arch} \\
Signal-consumption library & \textsc{architectural} & This paper,
Sec.~\ref{sec:arch} \\
Bittensor verification orchestration & \textsc{shadow only} (testnet,
57\,h, $5{,}097$ rounds) & Soak telemetry~\cite{omnirisk2026soak} \\
Validator-loss decomposition & \textsc{architectural} (deferred) &
Sec.~\ref{sec:vri:formal} \\
\bottomrule
\end{tabular}
\end{table*}

\subsection{Layer-1 Orchestration vs.\ Layer-2 Economic Verification}

A two-layer split discriminates what we have measured from what we
have not.

\paragraph*{Layer 1 (orchestration), measured.}
We have deployed an engineering shadow run of the OmniRisk validator
and miner pipeline against Bittensor testnet (netuid~60) since
2026-05-04~\cite{omnirisk2026soak}. Over an approximately 57-hour
window the run recorded \textbf{$5{,}097$ successful rounds} and
\textbf{zero authentication failures} across eleven autonomous Auth0
token refreshes (Table~\ref{tab:soak-telemetry}). The shadow run
validates the orchestration path: the canary layer can emit signed
signals to a Bittensor-style validation substrate continuously, under
autonomous credential rotation, without operator intervention.

\begin{table}[ht]
\centering
\small
\caption{Bittensor testnet shadow soak summary, netuid~60, 2026-05-04
to 2026-05-07. Reproduced from the engineering source-of-truth
artefact~\cite{omnirisk2026soak}. The table is the smallest possible
empirical surface for the layer-1 orchestration claim of
Section~\ref{sec:vri}; the layer-2 validator-loss decomposition is not
measured here.}
\label{tab:soak-telemetry}
\renewcommand{\arraystretch}{1.15}
\begin{tabular}{@{}lr@{}}
\toprule
\textbf{Metric} & \textbf{Value} \\
\midrule
Window (hours) & $\approx 57$ \\
Successful rounds & $5{,}097$ \\
Authentication failures & $0$ \\
Autonomous Auth0 token refreshes & $11$ \\
Operator interventions & $0$ \\
Validator-loss components reported & $0\,/\,3$ (deferred) \\
\bottomrule
\end{tabular}
\end{table}

\paragraph*{Layer 2 (economic verification), not yet.}
The three components of the validator loss
(Section~\ref{sec:vri:formal})---Brier accuracy, paired-query
inconsistency, binned miscalibration---are deferred. The smallest
experiment that would close the gap is to introduce a second miner
with a non-trivial protocol, accumulate $(c_i,y)$ pairs across
realised events, and report the three-component decomposition in a
follow-up. Until that runs, the verification claim is
\emph{architectural}, not measured. We state this honestly here because
the load-bearing claim of this paper is the canary-layer abstraction
plus its verifiability \emph{primitive}, not yet the empirical
validator-loss distribution at scale.

%% =====================================================================
\section{Use-Case Scenarios}\label{sec:uc}
%% =====================================================================

\subsection{Exploit Propagation}\label{sec:uc:exploit}

\paragraph*{Scenario.}
An attacker compromises a lending contract on chain~$A$ at time
$t_0$ and moves drained funds across a bridge to chain~$B$ at time
$t_0 + \delta$, where $\delta$ is bounded only by bridge latency and is
typically minutes. Chain-local analytics on chain~$A$ flag the exploit
at time $t_0 + \tau_A$ where $\tau_A$ is the chain-local detection
latency. In the absence of a canary layer, chain~$B$'s lending and
exchange-screening agents see the inbound wallet as clean for the
entire interval $[t_0+\delta,\,t_0+\tau_A+\delta_{B}]$, where
$\delta_{B}$ is the time it takes for the chain-$A$ analytics output to
appear, by some unstructured channel, on chain~$B$. The agent allows a
deposit, a swap, or a liquidation to settle, and the exploited capital
is laundered across chains.

\paragraph*{Canary-layer behaviour.}
With a canary layer, the wallet-level exploit signal generated on
chain~$A$ at $t_0+\tau_A$ is propagated under the typed schema and
arrives on chain~$B$ within the relayer-bound latency $\tau_R$.
Chain~$B$'s consumers reject the interaction or apply collateral
haircuts before settlement, provided $\tau_R \ll \delta_B$. The
operational metric of interest is the difference $\delta_B - \tau_R$
(Section~\ref{sec:eval:metrics}); the architectural metric is whether
chain~$B$'s consumer agents \emph{can} act on the propagated signal at
all under their constitutional action
surface~\cite{bai2022constitutional}---an agent whose policy whitelist
does not include canary-layer signal types is, by construction, blind
to the exploit.

\subsection{Credit Portability}\label{sec:uc:credit}

\paragraph*{Scenario.}
A wallet has a strong repayment and behavioural history on chain~$A$
(high repayment ratio, no liquidation events, sustained protocol
interaction over a multi-month window) but appears as a blank identity
on chain~$B$. Existing chain-specific reputation
systems~\cite{HasanBB23,AlmasoudHH20} cannot project the chain-$A$
history onto chain~$B$, so an under-collateralised lending protocol on
chain~$B$ either rejects the wallet or applies its conservative-default
haircut. The wallet's effective cost of capital is therefore
chain-local, even though its underlying creditworthiness is not.

\paragraph*{Canary-layer behaviour.}
With a canary layer carrying the wallet-level reputation signal under
a typed schema, an attached confidence, and a window of validity, the
lending agent on chain~$B$ extends risk-aware credit and reduces the
haircut by a function of the propagated score's confidence band. The
agentic-finance refinement is operationally important: a constitutional
Trader-role contract on chain~$B$ can declare a per-signal-type minimum
confidence, so a low-confidence cross-chain reputation signal is
silently downweighted rather than producing a deterministic credit
extension. RQ5 asks how far this generalises; we treat the use case
here as a structural argument, not yet a randomised experiment.

\subsection{Liquidity Stress}\label{sec:uc:liquidity}

\paragraph*{Scenario.}
A market-wide stress event manifests first as cross-chain withdrawal
correlation, bridge-flow asymmetry, and divergent quotes for the same
logical asset across L1 and L2 venues~\cite{gogol2024liquidity}, before
chain-local liquidity metrics on any single chain register the regime
change. Chain-local monitors therefore lag the systemic event by an
amount that, in the agentic regime, is sufficient for an autonomous
agent to drain liquidity from a pool whose state has already shifted.

\paragraph*{Canary-layer behaviour.}
A protocol-level signal (TVL contraction, oracle divergence) and a
market-level signal (correlated cross-chain withdrawal indicator) are
jointly consumed by liquidity-management agents whose constitutional
action surface includes pausing market-making, raising borrow rates
pre-emptively, or de-staking from at-risk pools. The architectural
prediction is that aggregating across chains shortens detection
latency relative to chain-local monitoring; the falsification condition
is stated in Section~\ref{sec:eval:metrics}.

%% =====================================================================
\section{Experimental Design and Evaluation}\label{sec:eval}
%% =====================================================================

This section sets out the experimental scaffold for the next-revision
evaluation. The architecture described in Sections~\ref{sec:arch}
and~\ref{sec:vri} is deployable; the empirical evaluation programme is
a road map of falsifiable measurements rather than a complete
empirical contribution at this revision.

\subsection{Datasets}

The evaluation draws on three dataset classes.
First, on-chain transaction streams across the eight chains covered by
the production deployment of the underlying signal engine
(Ethereum, BNB, Arbitrum, Optimism, Solana, Base, Polygon, Bittensor).
Second, wallet behaviour logs and labelled fraud cohorts from the
multichain DeFi fraud-detection
literature~\cite{palaiokrassas2023multichain,bhowmik2021,FanLZJB20}.
Third, protocol- and pool-level liquidity time series across L1 and
L2 venues, capturing the same logical asset across multiple
chains~\cite{gogol2024liquidity}.

\subsection{Metrics}\label{sec:eval:metrics}

We report three primary metrics.

\paragraph*{Detection latency}
The interval between the on-chain event that materialises a risk and
the moment the typed signal becomes consumable on a remote chain.
Hypothesis H1 predicts that cross-chain propagation reduces detection
latency relative to chain-local monitoring; the falsification
condition is that the propagated-signal latency is statistically
indistinguishable from the chain-local latency over a sufficiently
large event sample.

\paragraph*{Signal accuracy}
The directional accuracy and Brier calibration error of the signal
against realised outcomes. Calibration is reported alongside accuracy
because, in production deployments of the underlying prediction
engine, headline accuracy is class-imbalance-bounded and the survival
metric is the Brier
score~\cite{omnirisk2026predmetrics,omnirisk2026agentic}.

\paragraph*{Cross-chain coverage}
The fraction of qualifying events on each origin chain for which a
typed canary-layer signal is emitted within a fixed latency budget,
and the fraction of consuming chains that successfully decode and act
on the signal.

\subsection{Results: Orchestration Telemetry}

The only empirical surface in this revision is the orchestration soak
of Section~\ref{sec:vri}: $5{,}097$ successful rounds and zero
authentication failures over a 57-hour Bittensor testnet shadow
run~\cite{omnirisk2026soak}, with eleven autonomous Auth0 token
refreshes. The figure of merit here is operational: the orchestration
path holds under continuous unattended operation and meets the
preconditions for a multi-miner economic-verification experiment
without further operator scaffolding. We do not report validator-loss
distributions at this revision.

%% =====================================================================
\section{Discussion}\label{sec:disc}
%% =====================================================================

\subsection{Benefits}

The canary layer provides three categories of benefit. First,
\emph{earlier detection}: cross-chain propagation tightens the window
between exploit occurrence and consumer-side reaction
(Section~\ref{sec:uc:exploit}). Second, \emph{better protocol
decisions}: typed, confidence-bearing signals enable lending,
exchange-screening, and liquidity-management protocols to apply
context-aware policies rather than chain-local
defaults~\cite{HasanBB23}. Third, \emph{constitutional safety for
agents}: an autonomous agent under a constitutional
contract~\cite{bai2022constitutional} acts only on signals within its
declared type and confidence surface, so the canary layer narrows the
agent's effective action set to safer policies, especially under
liquidity stress and exploit-active regimes.

\subsection{Limitations}

We state the load-bearing limitations as falsification statements,
adopting the discipline of the companion focused
paper~\cite{magableh2026agenticiov}.

\paragraph*{Validator loss not yet measured at scale.}
The verification claim of Section~\ref{sec:vri} is supported only by
the orchestration soak; the three-component validator loss is not yet
measured against a live multi-miner metagraph. A follow-up that
introduces a second miner protocol and reports the decomposition
across realised events is the smallest experiment that would
falsify or substantiate the claim.

\paragraph*{Data-quality dependence.}
Wallet- and protocol-level signal quality is bounded by the quality of
upstream on-chain and off-chain ingestion. Coordinated wash trading
poisons the very signals the layer is meant to
propagate~\cite{cong2022crypto}; a sentiment-poisoned grey-literature
feed degrades the market-level signal class.

\paragraph*{False positives.}
A propagated risk signal that fires inappropriately is more harmful in
the agentic regime than in a human-trader regime, because an
autonomous agent acts on the signal at machine speed. False-positive
budgets must be declared per signal class, and the consumption layer
must support hysteresis and quorum requirements.

\paragraph*{Bridge-substrate trust.}
The canary layer inherits the trust assumptions of its underlying
transport. LayerZero's ultra-light-node design is itself a security
surface~\cite{ZhangZZL24,LiQXZZWX25,augusto2024sok}; an exploit of the
substrate is an exploit of the layer.

\subsection{Future Work}

Three lines of work follow directly. First, a multi-miner economic
verification of the validator loss against a Bittensor metagraph,
closing the layer-2 gap of Section~\ref{sec:vri}. Second, an
agentic consumer integration: a constitutional Trader-role policy
whose action surface is keyed to canary-layer signal types, building
on the policy-constrained execution
machinery~\cite{omnirisk2026agentic}. Third, predictive scoring under
distribution shift: AI-based scorers that adapt to evolving exploit
classes without sacrificing calibration on the cohorts they were
trained on.

%% =====================================================================
\section{Conclusion}\label{sec:concl}
%% =====================================================================

RayaChain introduces a missing infrastructure primitive for
decentralized agentic finance: an interoperability layer for
\emph{intelligence}, not value. Where existing protocols transport
value and arbitrary messages, the canary layer transports
evaluated risk and credit context under a typed schema with
confidence and a window of validity. We formalised the abstraction,
characterised the signal taxonomy, and demonstrated that the layer
becomes verifiable in expectation when paired with a Bittensor-style
decentralised validation subnet. The two empirical anchors are an
orchestration shadow soak ($5{,}097$ rounds, zero authentication
failures, eleven autonomous credential rotations) and a forward
agenda whose three falsifiable hypotheses are stated formally enough
that the next revision can refute them. The contribution is
architectural: a canary layer that lets autonomous agents on
heterogeneous chains act on the same evaluated context, rather than
on chain-local fragments of it.

%% =====================================================================
\section*{Author Contributions}
%% =====================================================================

We use the CRediT taxonomy. \textbf{Basel Magableh:}
conceptualisation, methodology, project administration, supervision,
writing -- original draft, writing -- review and editing.
\textbf{Ahod [surname TBD]:} contribution roles to be confirmed at the
time of submission. ORCIDs: Basel Magableh,
\texttt{0000-0003-2337-637X}; Ahod, \texttt{[ORCID TBD]}.

%% =====================================================================
\section*{Funding}
%% =====================================================================

This work received no external research funding. The on-chain shadow
soak telemetry referenced in Section~\ref{sec:vri} was collected on
Bittensor testnet (netuid~60) using infrastructure operated by the
Rayachain~Lab; no production funds were placed at risk during the
soak.

%% =====================================================================
\section*{Data Availability}
%% =====================================================================

The Bittensor shadow-soak telemetry summarised in
Table~\ref{tab:soak-telemetry} is committed alongside this manuscript
in the engineering source-of-truth
artefact~\cite{omnirisk2026soak}. The signal-schema definition of
Section~\ref{sec:arch} is reproducible from the typed-IDL listing in
that section; reference encoders and decoders are available to
reviewers under NDA on request.

%% =====================================================================
\section*{Conflict of Interest}
%% =====================================================================

The corresponding author is the operator of the OmniRisk production
deployment cited as the underlying signal-generation
engine~\cite{omnirisk2026agentic,omnirisk2026predmetrics}. No other
competing financial interests are disclosed.

%% =====================================================================
\section*{Ethics}
%% =====================================================================

No human-subjects research; no PII processing; no third-party data
acquisition under restricted licence. All on-chain data referenced are
publicly observable. The shadow-soak experiment of Section~\ref{sec:vri}
operated entirely on Bittensor testnet and did not act on any other
party's funds or data.

%% =====================================================================
%% Bibliography
%% =====================================================================

\bibliographystyle{IEEEtran}
\bibliography{references}

%% =====================================================================
%% Author biographies
%% =====================================================================

%% Author biographies (IEEEbiography blocks) deferred until Author~2
%% (Ahod) affiliation is confirmed by the editorial process.

\end{document}
